\section{The exact stride for cubic Jensen tests}
The preceding appendix bounds the number of tests when the polynomial degree is known. For continuous $V$ bounded below and of at most polynomial growth, consider testing just one infinite arithmetic progression of shifts. Persistence has the same pointwise-in-the-shift meaning as in Appendix~\ref{sec:finite-jensen}.
\begin{theorem}
If all shifts in $r+2\mathbb N_0$ are persistent for any integer $r\ge0$, then $V$ is a polynomial of degree at most two. For every stride $k\ge3$ and residue $r\bmod k$ there is a smooth strictly convex increasing nonpolynomial $V$ whose persistent shifts are exactly that residue class. Thus one parity, even with an arbitrary finite initial segment missing, detects every nonquadratic perturbation in this class.
\end{theorem}
\begin{proof}
The cubic discriminant expansion in the preceding appendix applies with the stated growth bounds. Persistence implies $\Delta^3m(n)=0$, where $m(n)=\mathbb EV(X_n)$ and $X_n$ has gamma shape $n+1/2$ and rate one.
Form a distribution supported on $[0,\infty)$ by
\[
 T(\phi)=-\int_0^\infty V(x)\frac{\dd^3}{\dd x^3}
       \bigl(x^{r+5/2}e^{-x}\phi(x)\bigr)\,\dd x .
\]
The least exponent at the origin is $r-1/2>-1$. Polynomial growth and exponential decay also let us test $T$ on polynomials and on exponentials $e^{zx}$ for $\Re z<1$. For $j\ge0$ and $p=r+j+1/2$, differentiating the integrand gives
\[
 T(x^j)=\Gamma(p+3)\Delta^3m(r+j).
\]
Indeed the four terms after differentiating have integrals proportional to $m(r+j+3),-3m(r+j+2),3m(r+j+1),-m(r+j)$ with the same factor. Put $F(z)=T(e^{zx})$. This function is analytic on $\Re z<1$ by dominated differentiation. If the progression of stride two is persistent, all even derivatives of $F$ at zero vanish. Hence $F(z)+F(-z)=0$ on $|\Re z|<1$ by analytic continuation. On the imaginary axis this says that $T$ plus its reflected distribution has zero Fourier transform. The distributions are tempered, and their Fourier transforms are the indicated functions (expand the derivative in the defining integral). Fourier uniqueness gives $T+T^{\rm refl}=0$. Their supports meet only at the origin, so $T$ vanishes on $(0,\infty)$. There $T=x^{r+5/2}e^{-x}V'''$ in distributions. The multiplier is smooth and strictly positive. Thus $V'''=0$, forcing a quadratic or lower polynomial, with the value at zero fixed by continuity. Appendix~\ref{sec:finite-jensen} supplies its converse.

For sharpness set $t=\tan(\pi/k)>0$, $h=(1-it)^{-1}$, with $\arg h=\pi/k$, and take
\[
 V(x)=Ax^2+Bx+\cos(tx+\theta).
\]
The gamma expectation of the oscillatory part is $\Re(e^{i\theta}h^{n+1/2})$. Its third difference is $\Re(e^{i\theta}h^{n+1/2}(h-1)^3)$. Choose
$\theta=\pi/2-(r+1/2)\pi/k-3\arg(h-1)$.
Since $h\ne1$, the third difference vanishes precisely at $n\equiv r\pmod k$. The quadratic and linear terms leave it unchanged. The oscillatory contribution to the second difference has absolute value at most $|h-1|^2$, since $|h|<1$. Choose $2A>\max(t^2,|h-1|^2)$ and $B>t$. Then $V'>0$, $V''>0$ and $\Delta^2m(n+1)>0$ everywhere. On the chosen progression the leading discriminant term is positive of order $\epsilon^3$. Off it the leading term is negative of order $\epsilon^2$. This gives the exact set of persistent shifts.
\end{proof}
This sampling statement concerns a perturbative test on a class of kernels; the smallness threshold for $\epsilon$ depends on the shift. The conductor-$47$ zero classification instead holds directly for the arithmetic functions and for every real parameter $c$.
